% Cette trace écrite à été développée par
% + Héloïse Troublé <heloise.trouble@ens-rennes.fr>
% + Erwan Tanguy-Legac <erwan.tanguy-legac@ens-rennes.fr>

% Elle est distribuée sous la licence Creative-Commons by-nc-sa 4.0
% (Attribution, Non-Commercial, Share Alike)
% qui peut être trouvée dans
% [le site de Creative Commons](https://creativecommons.org/licenses/by-nc-sa/4.0/)

\documentclass[12pt,a5paper]{article}
\usepackage[left=1.5cm, right=1.5cm, lines=45, top=0.3in, includehead]{geometry}
\usepackage[T1]{fontenc}
\usepackage[french]{babel}
\usepackage[utf8]{inputenc}
\usepackage{fancyhdr}
\usepackage{titling}
\usepackage{ragged2e}
\usepackage{subcaption}
\usepackage{wrapfig}
\usepackage{tikz}
\usetikzlibrary{quotes}

\DeclareCaptionFormat{custom}{#3}
\captionsetup{format=custom}

\setlength{\droptitle}{-2em}

\predate{}
\postdate{}
\date{\vspace{-12ex}}

\lhead{Informatique sans ordi}
\rhead{12 avril 2024}
\title{\vspace{-5ex}Journal du Pirate}

\pagenumbering{gobble}

\begin{document}

\maketitle
\thispagestyle{fancy}
\justifying

Des explorateurs ont trouvé le journal d’un grand pirate qui contient sans doute l’emplacement de son immense trésor. Vous devez les aider à savoir sur quelle île il a été caché. Le journal contient uniquement la description du chemin qu’a emprunté le pirate.
On a commencé par \textbf{modéliser} le chemin que le pirate a emprunté.

{
\centering
\begin{figure}[h]
  \centering
\begin{tikzpicture}[node distance=22.5mm, node/.style={thick, minimum size=16mm, rectangle, rounded corners=3mm, draw}]
  \node (init) [] {};
  \node (plage) [node, right of=init, node distance=12mm] {};
  \node (palmiers) [node, right of=plage] {};
  \node (cabane) [node, right of=palmiers] {};
  \node (volcan) [node, right of=cabane] {};
  \node (tresor) [node, right of=volcan] {};
  \draw [-stealth, thick] (init) to (plage);
  \draw (plage) edge["2h", thick, -stealth] (palmiers);
  \draw (palmiers) edge["3h", thick, -stealth] (cabane);
  \draw (cabane) edge["2h", thick, -stealth] (volcan);
  \draw (volcan) edge["1h", thick, -stealth] (tresor);
\end{tikzpicture}
\end{figure}
}

Après avoir modélisé le chemin du pirate, on a pu modéliser avec un schéma les différentes îles qui nous ont été distribuées pour vérifier
s'il était possible pour le pirate d'atteindre le trésor sur chaque île.
On a ainsi pu trouver la seule île sur lequel le trésor pouvait se cacher (en prenant en compte les temps de marche) !

\vspace{3mm}

Ce qu'on a fait ici ressemble fortement à un domaine de l'informatique qui est le diagnostic : on cherche à savoir si une machine fonctionne
normalement, c'est-à-dire s'il est possible ou non d'atteindre l'état \textit{en panne} depuis l'état de départ en suivant les transitions
possibles.

\begin{center}
\begin{tikzpicture}[text height=1.5ex, text depth=.25ex, node/.style={thick, minimum size=8mm, rectangle, draw}]
  \node (empty) at (0, 1) [] {};
  \node (init) at (0, 0) [node] {début};
  \node (marche) at (-0.2, -1.1) [node] {marche};
  \node (veille) at (1.7, -1.1) [node] {veille};
  \node (arret) at (0, -2.2) [node] {fin de tâche};
  \node (panne) at (-2.2, -1.1) [node] {panne};

  \draw [-stealth, thick] (empty) to (init);
  \draw [-stealth, thick] (init) to (marche);
  \draw [stealth-stealth, thick] (marche) to (veille);
  \draw [-stealth, thick] (marche) to (arret);
  \draw [-stealth, thick] (arret) to (veille);
  \draw [-stealth, thick] (panne) to (init);
  \draw [-stealth, thick] (arret) to (panne);
\end{tikzpicture}
\end{center}

Dans l'exemple ci-dessus, où on a modélisé une machine, on peut s'assurer qu'il est impossible pour la machine de tomber en panne avant d'avoir
fini au moins une fois sa tâche.
\end{document}
